Se sabe que un problema espacial tiene la ecuación
\[
-X''=\lambda X,\qquad 0<x<L,
\]
con \(L>0\) conocido. En cada extremo se impone exactamente una de estas dos condiciones: valor nulo \(X=0\) o derivada nula \(X'=0\). Se desconoce cuál se usa en cada extremo.

Un registro sintético exacto contiene los dos menores autovalores \emph{estrictamente positivos}:
\[
\lambda_1^+=\left(\frac{\pi}{L}\right)^2,
\qquad \lambda_2^+=\left(\frac{2\pi}{L}\right)^2.
\]
El registro no ha comprobado si \(\lambda=0\) es autovalor. Determine todas las parejas de condiciones compatibles con la información disponible. ¿Bastaría comprobar la existencia de una autofunción de autovalor cero para identificar la pareja? Justifique ambas decisiones.
